\subsection{Witnesses for the actuarial axioms}\label{supp:actwit}

\begin{table}[ht]
\centering\small
\begin{tabularx}{\textwidth}{@{}l>{\raggedright\arraybackslash}p{4.6cm}>{\raggedright\arraybackslash}X@{}}
\toprule
\textbf{Dropped} & \textbf{Witness on $\Xp$} & \textbf{Reason}\\
\midrule
(A2) & $\E_Q/\E$ & $N_0$ lies in a closed hyperplane, so (A0) holds; not law invariant.\\
(A3) & $C_h/\E$, $h=t+\tfrac14\sin2\pi t$ & $\rho(\1_A)>\rho(\1_B)$ for $\1_A\le\1_B$, $P(A)=0.4,P(B)=0.6$.\\
(A3$_s$) & $C_h/\E$, $h$ continuous non-BV & Extensions exist; $h$ is not of bounded variation.\\
(A4a) & $\E/\max(C_{\sqrt t},C_{\min(1,1.5t)})$ & Premium is not comonotonic additive (numbers above); (A4b) holds.\\
(A4b) & $\max(C_{\sqrt t},C_{\min(1,1.5t)})/\E$ & $\bar\Pi=\E$ additive, but $\bar\Pi\LR_0$ is not.\\
(A5) & $C_h/C_g$ with $g=(2t-1)^+$, $h=g^2$, and $\LR_0:=1$ where $C_g=0$ & $\LR_0$ jumps: $\LR_0(\1_B+n^{-1}\1_C)=\tfrac12$ while $\LR_0(\1_B)=1$ for $P(B)\le\tfrac12$, $C\supset B$, $P(C)=\tfrac34$.\\
\bottomrule
\end{tabularx}
\caption{Witnesses for (A2)--(A5) on the actuarial domain. In each case all other axioms hold.}\label{ex:A5}
\end{table}